% 5. Rezultati
\section{Rezultati}
\label{sec:rezultati}

Tabele~\ref{tab:pra_results} i~\ref{tab:gra_results} sumiraju rezultate za PRA odnosno GRA režim kroz tri pokazatelja: PCC kao meru linearne povezanosti sa ljudskim ocenama, odnos prosečnih ocena (prosečna ocena sistema / prosečna ocena nastavnika) kao pokazatelj sistematskog odstupanja (blagosti ili strogosti) i odnos RMSE/STDDEV kao meru relativnog odstupanja od varijabilnosti ljudskog ocenjivanja.
Sve vrednosti izračunate su nad objedinjenim skupom svih ocenjenih odgovora, a ne kao prosek vrednosti po pojedinačnim skupovima podataka.

Svaka konfiguracija (model $\times$ strogost $\times$ režim) pokrenuta je jednom, pa rezultati odražavaju jedno uzorkovanje iz svakog modela.
Temperatura se u glavnoj mreži konfiguracija ne varira:
modeli porodice GPT rade na podrazumevanoj vrednosti svog provajdera,
a modelima DeepSeek prosleđuje se vrednost $0$ (odeljak~\ref{sec:pipeline}).
Njen uticaj ispitan je odvojeno, na podskupu odgovora, u odeljku~\ref{sec:rezultati-temperatura}.

\subsection{Performanse modela}
\label{sec:performanse}

\begin{table}[ht]
\centering
\caption{Metrike performansi za PRA režim kroz modele i konfiguracije strogosti.
Najbolji rezultati za svaku metriku i konfiguraciju strogosti označeni su podebljano.
Za odnos prosečnih ocena najboljom se smatra vrednost najbliža 1{,}0.}
\label{tab:pra_results}
\resizebox{\textwidth}{!}{%
\begin{tabular}{l ccc ccc ccc}
\toprule
\textbf{Model} &
\multicolumn{3}{c}{\textbf{PCC}} &
\multicolumn{3}{c}{\textbf{Odnos proseka}} &
\multicolumn{3}{c}{\textbf{RMSE/STDDEV}} \\
\cmidrule(lr){2-4}
\cmidrule(lr){5-7}
\cmidrule(lr){8-10}
& blag & neutr. & strog
& blag & neutr. & strog
& blag & neutr. & strog \\
\midrule
GPT-4o-mini       & 0{,}82 & 0{,}80 & 0{,}77 & \textbf{0{,}96} & 0{,}93 & \textbf{0{,}97} & 0{,}57 & 0{,}61 & 0{,}65 \\
GPT-4o            & 0{,}80 & 0{,}84 & 0{,}79 & 0{,}90 & \textbf{1{,}00} & 0{,}79 & 0{,}65 & 0{,}55 & 0{,}74 \\
GPT-5             & \textbf{0{,}89} & \textbf{0{,}88} & 0{,}77 & 0{,}94 & 0{,}91 & \textbf{1{,}00} & \textbf{0{,}47} & \textbf{0{,}50} & 0{,}65 \\
DeepSeek-Chat     & 0{,}81 & 0{,}82 & \textbf{0{,}83} & 0{,}80 & 0{,}84 & 0{,}86 & 0{,}71 & 0{,}65 & \textbf{0{,}61} \\
DeepSeek-Reasoner & 0{,}87 & 0{,}86 & 0{,}78 & 0{,}88 & 0{,}86 & 0{,}95 & 0{,}53 & 0{,}56 & 0{,}63 \\
\bottomrule
\end{tabular}
}
\end{table}

Vrednosti u tabeli~\ref{tab:pra_results} izračunate su nad svih 583 odgovora, osim u dva slučaja u kojima ocenjivanje nije uspelo za sve odgovore: neutralna konfiguracija modela GPT-4o obuhvata 399 odgovora, a neutralna konfiguracija modela DeepSeek-Reasoner 555 odgovora.\footnote{Nedostajuće ocene posledica su prekinutog pokretanja, odnosno trajno neuspešnih API poziva, koji se prema opisu iz odeljka~\ref{sec:pipeline} označavaju i isključuju iz analize.}

Modeli usmereni na rezonovanje pokazuju najjače ukupne performanse.
GPT-5 postiže najvišu korelaciju u blagoj i neutralnoj konfiguraciji (PCC $=0{,}89$ i $0{,}88$) uz najniže odnose RMSE/STDDEV ($0{,}47$ i $0{,}50$), a DeepSeek-Reasoner ga prati u stopu ($0{,}87$ i $0{,}86$).
GPT-4o je konkurentan u neutralnoj konfiguraciji (PCC $=0{,}84$), gde postiže i odnos prosečnih ocena praktično jednak $1{,}0$, dok su GPT-4o-mini i DeepSeek-Chat dosledno nešto slabiji.

Stroga konfiguracija menja poredak: u njoj najvišu korelaciju postiže DeepSeek-Chat (PCC $=0{,}83$), a GPT-5 i GPT-4o-mini padaju na $0{,}77$.
DeepSeek-Chat je pri tom jedini model kod koga korelacija monotono raste sa strogošću ($0{,}81 \rightarrow 0{,}82 \rightarrow 0{,}83$), dok kod svih ostalih stroga konfiguracija daje najniže vrednosti.
Kod modela GPT-4o stroga konfiguracija istovremeno unosi izraženo sistematsko odstupanje naniže (odnos proseka $0{,}79$), što znači da model ocenjuje znatno strože od nastavnika.

Odnosi prosečnih ocena pokazuju da modeli iz porodice GPT ostaju bliži prosecima ljudskog ocenjivanja od modela DeepSeek.
DeepSeek-Chat dosledno ocenjuje konzervativnije od nastavnika u svim konfiguracijama (odnosi $0{,}80$--$0{,}86$), dok GPT-4o-mini ostaje u rasponu $0{,}93$--$0{,}97$.

\begin{table}[ht]
\centering
\caption{Metrike performansi za GRA režim.
Najbolji rezultati za svaku metriku i konfiguraciju strogosti označeni su podebljano.}
\label{tab:gra_results}
\resizebox{\textwidth}{!}{%
\begin{tabular}{l ccc ccc ccc}
\toprule
\textbf{Model} &
\multicolumn{3}{c}{\textbf{PCC}} &
\multicolumn{3}{c}{\textbf{Odnos proseka}} &
\multicolumn{3}{c}{\textbf{RMSE/STDDEV}} \\
\cmidrule(lr){2-4}
\cmidrule(lr){5-7}
\cmidrule(lr){8-10}
& blag & neutr. & strog
& blag & neutr. & strog
& blag & neutr. & strog \\
\midrule
GPT-4o-mini       & 0{,}63 & 0{,}65 & 0{,}61 & 0{,}79 & 0{,}77 & 0{,}82 & 0{,}85 & 0{,}84 & 0{,}84 \\
GPT-4o            & 0{,}69 & 0{,}75 & 0{,}65 & 0{,}78 & \textbf{0{,}88} & 0{,}61 & 0{,}86 & 0{,}69 & 1{,}04 \\
GPT-5             & \textbf{0{,}85} & \textbf{0{,}84} & \textbf{0{,}73} & \textbf{0{,}87} & 0{,}83 & \textbf{0{,}89} & \textbf{0{,}57} & \textbf{0{,}61} & \textbf{0{,}73} \\
DeepSeek-Chat     & 0{,}43 & 0{,}50 & 0{,}47 & 0{,}37 & 0{,}52 & 0{,}42 & 1{,}42 & 1{,}22 & 1{,}34 \\
DeepSeek-Reasoner & 0{,}79 & 0{,}79 & 0{,}70 & 0{,}81 & 0{,}77 & 0{,}87 & 0{,}70 & 0{,}72 & 0{,}77 \\
\bottomrule
\end{tabular}
}
\end{table}

Tabela~\ref{tab:gra_results} prikazuje rezultate za \textbf{GRA} režim, u kome referentne odgovore automatski generiše jezički model umesto nastavnika.
Očekivano, rezultati u GRA režimu su niži od onih dobijenih sa referentnim odgovorima nastavnika.
Ovaj pad performansi odražava dodatnu neizvesnost koju unose automatski generisani referentni odgovori, koji se mogu razlikovati po strukturi, formulaciji ili nivou detalja od rešenja koje je nastavnik imao na umu.

Modeli za rezonovanje najbolje podnose prelazak u GRA režim.
GPT-5 zadržava snažno slaganje (PCC $=0{,}85$ u blagoj i $0{,}84$ u neutralnoj konfiguraciji) i jedini u svim konfiguracijama drži RMSE/STDDEV ispod $0{,}75$, a DeepSeek-Reasoner ostaje na $0{,}79$.
To sugeriše da su sposobniji modeli za rezonovanje istovremeno bolji i u generisanju prikladnih referentnih odgovora i u ocenjivanju u odnosu na njih.

DeepSeek-Chat je u GRA režimu izrazito neupotrebljiv: korelacija pada na $0{,}43$--$0{,}50$, odnos prosečnih ocena na $0{,}37$--$0{,}52$, a RMSE/STDDEV prelazi $1{,}2$ u svim konfiguracijama, dakle njegove ocene odstupaju od ljudskih više nego što ljudske ocene variraju međusobno.
GPT-4o pokazuje isti problem u strogoj konfiguraciji (RMSE/STDDEV $=1{,}04$, odnos proseka $0{,}61$).

\subsection{Poređenje PRA i GRA režima kroz modele i nivoe strogosti}
\label{sec:pra-vs-gra}

Kroz sve evaluirane modele i konfiguracije strogosti, PRA režim dosledno nadmašuje GRA režim, i po korelaciji sa ljudskim ocenama i po normalizovanim odnosima greške.

\subsubsection{Korelacija sa ljudskim ocenjivanjem}

Prelaskom u GRA režim korelacije opadaju kod svih modela, ali neujednačeno.
Kod modela GPT-5 pad je najmanji (sa $0{,}88$ na $0{,}84$ u neutralnoj konfiguraciji), a kod modela DeepSeek-Chat najveći (sa $0{,}82$ na $0{,}50$).
Ova razlika pokazuje da kvalitet automatski generisanog referentnog odgovora, a ne samo sposobnost ocenjivanja, presudno određuje performanse u GRA režimu.

\subsubsection{Uticaj konfiguracije strogosti}

Ne postoji konfiguracija strogosti koja je najbolja za sve modele, i optimalna konfiguracija se razlikuje između dva režima.
U PRA režimu najvišu korelaciju blaga konfiguracija daje kod modela GPT-5, GPT-4o-mini i DeepSeek-Reasoner, neutralna kod modela GPT-4o, a stroga kod modela DeepSeek-Chat.
U GRA režimu neutralna konfiguracija je najbolja kod svih modela osim GPT-5, kod koga je najbolja blaga; kod modela DeepSeek-Reasoner razlika između blage i neutralne konfiguracije zanemarljiva je (oba $0{,}79$).

Stroga konfiguracija je u oba režima najlošija za četiri od pet modela, i to ne samo po korelaciji nego i po sistematskom odstupanju, kao kod modela GPT-4o u GRA režimu.

\subsubsection{Sistematsko odstupanje i odnosi greške}

Dok u PRA režimu odnosi prosečnih ocena ostaju blizu $1{,}0$ (odeljak~\ref{sec:performanse}), u GRA režimu svi modeli ocenjuju niže od nastavnika, što je u skladu sa nalazom iz odeljka~\ref{sec:cross-dataset} da su generisani referentni odgovori po pravilu detaljniji od nastavničkih, pa studentski odgovori u odnosu na njih izgledaju nepotpuniji.

Isto razdvajanje pokazuju i odnosi greške: u PRA režimu svi modeli ostaju ispod vrednosti $0{,}75$, dok se u GRA režimu granica $1{,}0$ prelazi u dva slučaja, kod modela DeepSeek-Chat u svim konfiguracijama ($1{,}22$--$1{,}42$) i kod modela GPT-4o u strogoj konfiguraciji ($1{,}04$).

\subsubsection{Uporedni pregled performansi}

Kroz metrike korelacije, sistematskog odstupanja i greške izdvajaju se tri zaključka:

\begin{itemize}
    \item Referentni odgovori nastavnika (PRA) značajno poboljšavaju pouzdanost ocenjivanja,
    kod svih modela i svih konfiguracija strogosti.
    \item Modeli usmereni na rezonovanje najbolje se slažu sa ljudskim ocenjivačima:
    GPT-5 je ukupno najjači, uz najniže odnose greške,
    a DeepSeek-Reasoner pokazuje da ta prednost nije vezana za jednog provajdera.
    Izraženija je u GRA režimu, gde model mora sam da konstruiše referentni standard.
    \item Standardni modeli su upotrebljivi, ali osetljiviji na postavku:
    GPT-4o je konkurentan u neutralnoj konfiguraciji, dok mu stroga u GRA režimu narušava upotrebljivost,
    GPT-4o-mini nudi povoljan odnos cene i performansi u PRA režimu,
    a DeepSeek-Chat je upotrebljiv samo u PRA režimu.
\end{itemize}

\subsection{Analiza PCC-a po skupovima podataka}
\label{sec:cross-dataset}

Tabele~\ref{tab:pcc-reference} i~\ref{tab:pcc-generated} prikazuju vrednosti PCC-a po kursevima za PRA odnosno GRA ocenjivanje, u neutralnoj konfiguraciji strogosti za svaki model.
Budući da se broj odgovora po kursevima znatno razlikuje (tabela~\ref{tab:dataset}), vrednosti za manje kurseve treba tumačiti sa rezervom: kod kurseva DDS i HTML5 korelacija je izračunata nad po deset odgovora.

Primetna razlika između PRA i GRA rezultata javlja se na kursevima gde se očekuje veoma kratak i precizan studentski odgovor, kao što su WP i TSW.
Na ovim skupovima podataka vrednosti PCC-a za GRA opadaju u odnosu na PRA (tabela~\ref{tab:pcc-reference} naspram tabele~\ref{tab:pcc-generated}); kod kursa TSW pad je najizraženiji, sa $0{,}81$--$0{,}93$ na $0{,}46$--$0{,}72$.
Jedan od razloga je taj što automatski generisani referentni odgovori teže da budu znatno duži i opisniji od sažetih rešenja koja daje nastavnik.
Kao posledica, model često dodeljuje visoke ocene odgovorima koji ispravno opisuju traženi postupak, ali ga zapravo ne izvršavaju.
Nasuprot tome, nastavnik je takve odgovore ocenjivao strogo i po pravilu dodeljivao nula poena kada tražena radnja nije sprovedena.
Ovaj raskorak između opisne i proceduralne tačnosti smanjuje slaganje sa ljudskim ocenjivanjem.

Ovakvo ponašanje može se ublažiti uključivanjem kratke napomene nastavnika koja eksplicitno naglašava da li je tačno objašnjenje samo po sebi dovoljno ili student mora u potpunosti da izvrši traženi zadatak.
Takvo usmerenje pomaže modelu da kriterijume ocenjivanja bolje uskladi sa očekivanjima nastavnika.

Kod skupova podataka gde se očekuju duži, opisniji odgovori (npr.\ DDS i HTML5), razlike između PRA i GRA su znatno manje.
U tim slučajevima generisani referentni odgovori po strukturi i nivou detalja liče na odgovore nastavnika, što vodi sličnom ponašanju pri ocenjivanju i uporedivim vrednostima PCC-a u oba režima.
Na kursu HTML5 vrednosti su praktično identične u oba režima ($0{,}93$--$0{,}97$).

Najveći kurs, Programiranje 2, sa 477 od ukupno 583 odgovora, ponaša se u skladu sa objedinjenim rezultatima: GPT-5 je najbolji u oba režima ($0{,}88$ i $0{,}84$), a DeepSeek-Chat pokazuje najveći pad između režima (sa $0{,}80$ na $0{,}42$).

\begin{table}[ht]
\centering
\caption{PCC po kursevima za PRA režim (neutralna konfiguracija strogosti).
U zagradi je broj odgovora nad kojim je vrednost izračunata.}
\label{tab:pcc-reference}
\resizebox{\columnwidth}{!}{%
\begin{tabular}{lccccc}
\toprule
\textbf{Kurs} & \textbf{DeepSeek-Chat} & \textbf{DeepSeek-Reasoner} & \textbf{GPT-4o} & \textbf{GPT-4o-mini} & \textbf{GPT-5} \\
\midrule
Prog 2 & 0{,}80 (477) & 0{,}85 (449) & 0{,}82 (293) & 0{,}79 (477) & \textbf{0{,}88} (477) \\
TSW    & 0{,}81 (40) & 0{,}89 (40) & \textbf{0{,}93} (40) & 0{,}86 (40) & 0{,}91 (40) \\
WP     & 0{,}78 (46) & 0{,}74 (46) & \textbf{0{,}81} (46) & 0{,}78 (46) & 0{,}75 (46) \\
DDS    & 0{,}80 (10) & 0{,}82 (10) & 0{,}61 (10) & 0{,}86 (10) & \textbf{0{,}92} (10) \\
HTML5  & 0{,}96 (10) & 0{,}94 (10) & \textbf{0{,}97} (10) & \textbf{0{,}97} (10) & 0{,}94 (10) \\
\bottomrule
\end{tabular}
}
\end{table}

\begin{table}[ht]
\centering
\caption{PCC po kursevima za GRA režim (neutralna konfiguracija strogosti).
U zagradi je broj odgovora nad kojim je vrednost izračunata.}
\label{tab:pcc-generated}
\resizebox{\columnwidth}{!}{%
\begin{tabular}{lccccc}
\toprule
\textbf{Kurs} & \textbf{DeepSeek-Chat} & \textbf{DeepSeek-Reasoner} & \textbf{GPT-4o} & \textbf{GPT-4o-mini} & \textbf{GPT-5} \\
\midrule
Prog 2 & 0{,}42 (477) & 0{,}81 (477) & 0{,}74 (477) & 0{,}64 (477) & \textbf{0{,}84} (477) \\
TSW    & 0{,}47 (40) & 0{,}46 (40) & \textbf{0{,}72} (40) & 0{,}62 (40) & 0{,}65 (40) \\
WP     & 0{,}75 (46) & 0{,}66 (46) & \textbf{0{,}77} (46) & 0{,}56 (46) & 0{,}71 (46) \\
DDS    & 0{,}84 (10) & 0{,}86 (10) & 0{,}60 (10) & 0{,}51 (10) & \textbf{0{,}86} (10) \\
HTML5  & 0{,}96 (10) & 0{,}94 (10) & 0{,}96 (10) & \textbf{0{,}97} (10) & 0{,}93 (10) \\
\bottomrule
\end{tabular}
}
\end{table}

\subsection{Dodatni eksperimenti}
\label{sec:dodatni}

Pored glavne mreže konfiguracija sprovedena su i dva manja eksperimenta,
koji ispituju dve odluke koje ta mreža ne pokriva:
da li se posrednički referentni odgovor može izostaviti i umesto njega u prompt umetnuti ceo nastavni materijal,
i da li variranje temperature menja slaganje sa ljudskim ocenjivanjem.
Zbog ograničenog obuhvata prikazuju se odvojeno od glavnih rezultata;
ni jedan od njih ne menja nalaze iz prethodnih odeljaka.

\subsubsection{Ocenjivanje iz nastavnog materijala}
\label{sec:rezultati-udzbenik}

Treći režim ocenjivanja, opisan u odeljku~\ref{sec:rezim-udzbenik},
ne koristi referentni odgovor, već se ceo tekst udžbenika umeće u svaki prompt za ocenjivanje.
Zbog cene je pokrenut samo za modele GPT-4o-mini i GPT-4o,
pa se nalazi odnose isključivo na ta dva modela.
Tabela~\ref{tab:udzbenik} prikazuje njegove metrike,
uz vrednosti PCC-a u PRA i GRA režimu kao referencu.

\begin{table}[ht]
\centering
\caption{Metrike režima sa nastavnim materijalom u promptu,
uz vrednosti PCC-a u PRA i GRA režimu kao referencu.
Za model GPT-4o-mini sve vrednosti izračunate su nad svih 583 odgovora.
Za model GPT-4o režim sa nastavnim materijalom obuhvata 581 odgovor u neutralnoj i 454 u strogoj konfiguraciji,
dok je blaga konfiguracija ostala praktično nepokrenuta (12 odgovora) pa se ne prikazuje.
Referentna PRA vrednost za GPT-4o u neutralnoj konfiguraciji izračunata je nad 399 odgovora,
kao u tabeli~\ref{tab:pra_results}.}
\label{tab:udzbenik}
\resizebox{\columnwidth}{!}{%
\begin{tabular}{ll ccc cc}
\toprule
\multirow{2}{*}{\textbf{Model}} & \multirow{2}{*}{\textbf{Strogost}} &
\multicolumn{3}{c}{\textbf{PCC}} &
\multicolumn{2}{c}{\textbf{Režim sa materijalom}} \\
\cmidrule(lr){3-5}
\cmidrule(lr){6-7}
& & PRA & materijal & GRA & odnos proseka & RMSE/STDDEV \\
\midrule
GPT-4o-mini & blaga     & 0{,}82 & 0{,}70 & 0{,}63 & 0{,}92 & 0{,}73 \\
GPT-4o-mini & neutralna & 0{,}80 & 0{,}70 & 0{,}65 & 0{,}92 & 0{,}73 \\
GPT-4o-mini & stroga    & 0{,}77 & 0{,}63 & 0{,}61 & 0{,}89 & 0{,}80 \\
\midrule
GPT-4o      & neutralna & 0{,}84 & 0{,}67 & 0{,}75 & 0{,}86 & 0{,}79 \\
GPT-4o      & stroga    & 0{,}79 & 0{,}62 & 0{,}65 & 0{,}68 & 0{,}97 \\
\bottomrule
\end{tabular}
}
\end{table}

Kod oba modela režim sa nastavnim materijalom ostaje jasno ispod PRA režima,
sa vrednostima PCC-a u rasponu $0{,}62$--$0{,}70$ naspram $0{,}77$--$0{,}84$.
U odnosu na GRA režim nema dosledne prednosti:
kod modela GPT-4o-mini materijal u promptu daje nešto bolje slaganje u sve tri konfiguracije
($0{,}70$ naspram $0{,}65$ u neutralnoj),
a kod modela GPT-4o osetno slabije ($0{,}67$ naspram $0{,}75$).
Izostavljanje posredničkog referentnog odgovora, dakle,
ne donosi sistematsko poboljšanje kod pokrivenih modela.

Stroga konfiguracija se i u ovom režimu ponaša kao u GRA režimu:
kod modela GPT-4o povlači odnos prosečnih ocena na $0{,}68$, a odnos RMSE/STDDEV na granicu upotrebljivosti ($0{,}97$).

Pitanje iz odeljka~\ref{sec:rezim-udzbenik} time je delimično odgovoreno:
prompt je približno 55 puta veći nego u PRA režimu, a slaganje nije bolje,
pa dodatni kontekst ne opravdava svoju cenu.
Kod modela usmerenih na rezonovanje ishod bi mogao biti drugačiji,
budući da oni najbolje podnose prelazak u GRA režim,
ali u ovom režimu nisu pokretani.

\subsubsection{Uticaj temperature}
\label{sec:rezultati-temperatura}

Temperatura je parametar koji upravlja slučajnošću pri uzorkovanju odgovora modela:
niže vrednosti daju predvidljiviji, a više raznovrsniji izlaz.
Kako se u glavnoj mreži konfiguracija ne postavlja eksplicitno,
posebno je ispitano da li njeno variranje menja slaganje sa nastavnicima.
Eksperiment je pokrenut za model GPT-4o u neutralnoj konfiguraciji i PRA režimu,
sa temperaturama $0{,}5$ i $1{,}0$,
nad 305 odgovora: 285 odgovora sa dva od tri roka kursa Programiranje 2
i po deset odgovora sa kurseva DDS i HTML5.
Rezultati su prikazani u tabeli~\ref{tab:temperatura}.

\begin{table}[ht]
\centering
\caption{Uticaj temperature na metrike ocenjivanja
(model GPT-4o, neutralna konfiguracija strogosti, PRA režim, 305 odgovora).}
\label{tab:temperatura}
\begin{tabular}{lccc}
\toprule
\textbf{Temperatura} & \textbf{PCC} & \textbf{Odnos proseka} & \textbf{RMSE/STDDEV} \\
\midrule
podrazumevana & 0{,}82 & 1{,}02 & 0{,}58 \\
$0{,}5$       & 0{,}83 & 1{,}02 & 0{,}57 \\
$1{,}0$       & 0{,}82 & 1{,}00 & 0{,}58 \\
\bottomrule
\end{tabular}
\end{table}

Sve tri metrike poklapaju se do druge decimale,
pa variranje temperature u rasponu od $0{,}5$ do $1{,}0$ ne menja slaganje sa nastavnicima na merljiv način.
Pri tumačenju ovog nalaza treba imati u vidu da red sa podrazumevanom temperaturom
dolazi iz zasebnog pokretanja, a da podrazumevana vrednost koju provajder primenjuje
na model GPT-4o iznosi $1{,}0$.
Prvi i treći red tabele su, dakle, na istoj nominalnoj temperaturi,
pa razlika između njih ne meri uticaj temperature nego variranje između dva pokretanja istog modela,
i po veličini je uporediva sa svim ostalim razlikama u tabeli.

Jedini izraženiji pomak javlja se na kursu DDS, gde PCC raste sa $0{,}61$ na $0{,}89$ i $0{,}83$.
I on je, u skladu sa rezervom iz odeljka~\ref{sec:cross-dataset},
posledica pojedinačnih ocena među samo deset odgovora, a ne temperature.
Zbog ograničenog obuhvata ovaj nalaz treba čitati kao proveru da izbor temperature
nije prikriveni izvor razlika u glavnim rezultatima,
a ne kao potpuno ispitivanje njenog uticaja.
